\runninghead{Thoughts of Peace, Cries from the Depths}
\properlane{cry-from-captivity}{Thoughts of Peace for Those Who Cry from the Depths: Prayer Commanded within a Promise}

In the Introit God speaks through Jeremias to the exiles in Babylon, and what
he says is a promise: he thinks thoughts of peace, he will hear, and
\latin{redúcam captivitátem vestram}, ``I will bring back your captivity.'' The same
antiphon tells them \latin{invocábitis me}, ``you shall call upon me,'' and
later in the Mass the people do call: \latin{De profúndis clamávi}, ``from the
depths I have cried,'' once in the Alleluia and twice in the Offertory. If God
has already said what he will do, why must he be asked? St Jerome answers that
question on the verses of Jeremias the Introit takes. St Hilary of Poitiers,
on the psalm the Mass sings at the Alleluia and the Offertory, asks a
question of his own, and his answer bears on Jerome's. St Augustine says who
the captives are and what their captivity is: the captivity of sin, from
which Christ sets them free.

\subsection{What the Introit keeps of Jeremias's letter}
\label{sec:captivity-introit}

The Introit takes three pieces of Jeremias's letter and puts them together. Its
words, God's own, keep the thoughts of peace, the calling and the hearing,
and the bringing back of the captivity. Between them the Missal leaves out
four things that the letter says: that God thinks peace ``to give you an end
and patience''; that the exiles ``shall go: and \dots\ shall pray''; that
they will seek and find him ``when you shall seek me with all your heart,''
and that he ``will be found'' by them (vv.~13--14); and that he will ``gather
you out of all nations.'' What remains is God's intention, the people's call,
God's hearing and the return, in that order.

In the letter the exiles are told to call on God while they are still in
Babylon, after the return has been promised and before it comes; it will come
only when ``the seventy years shall begin to be accomplished in Babylon''
(Jer 29:10). The householder of the Gospel also tells his servants beforehand
what he will do and when: \latin{in témpore messis dicam messóribus}, ``in the
time of the harvest I will say to the reapers,'' and until then both are to
grow. In the parable as in the letter, the one who promises fixes the term
and tells those who wait what to do until it comes.

\subsection{St Jerome: God could do it without their prayer}
\label{sec:captivity-jerome}

St Jerome's commentary on the letter begins from the false prophets. They
claimed to know what God would do; Jeremias answers with \latin{Ego enim scio
cogitationes quas ego cogito}, ``I know the thoughts that I think,'' and
Jerome draws the consequence: \latin{Futurorum igitur scientia soli Deo
competit}, ``knowledge of things to come therefore belongs to God alone.''
The thoughts of peace are God's own, known to him before anyone prays. The
\latin{finem et patientiam} that God says he will give are \latin{finem
captivitatis, et patientiam laborum praesentium, sive spem futurorum}, the end
of the captivity, and patience in present labours, or hope of what is to come
(\work{In Hieremiam} V, on 29:10--14, PL~24 col.~893). The Introit leaves
these words out. The Epistle asks the Colossians to put on
\latin{patiéntiam}, and earlier in the chapter Paul has told them that their
life ``is hid with Christ in God'' and that ``when Christ shall appear, who is
your life, then you also shall appear with him in glory'' (Col 3:3--4). They
too hold a promise not yet fulfilled, so Jerome's ``patience in present
labours, or hope of what is to come'' describes their patience as well as the
exiles', and the patience the antiphon leaves unsung is asked for in the
Epistle.

On the next verse, ``you shall call upon me \dots\ and I will hear you,''
Jerome writes: \latin{Certe
absque invocatione et oratione captivorum, poterat implere Dominus quod
promiserat; sed hortatur eos ad preces, ut quod promissum est, mereantur
accipere}. ``Certainly the Lord could have fulfilled what he had promised
without the invocation and prayer of the captives; but he urges them to pray,
so that they may be worthy to receive what has been promised.'' On verse~13,
``you shall seek me, and shall find me,'' he sets the Lord's words in the
Gospel: \latin{Petite, et accipietis}, ``Ask, and you shall receive'' (Mt
7:7).

Jerome also reads the letter anagogically: \latin{tamdiu in saeculi hujus
confusione versamur, quamdiu requiem septenarii numeri mereamur accipere},
``we remain in the confusion of this age until we are found worthy to receive
the rest of the number seven''; then, repentance having been received, God
fulfils what he promised, \latin{et reducamur ad locum nostrum Ecclesiam},
``and we are brought back to our own place, the Church.''

\subsection{St Hilary and St John Chrysostom on the Mass's own psalm}
\label{sec:captivity-psalm}

The Alleluia and the Offertory both sing \latin{De profúndis clamávi ad te,
Dómine: Dómine, exáudi oratiónem meam}. St Hilary of Poitiers, expounding the
psalm, asks whether ``hear my voice'' and the next words, ``let thy ears be
attentive to the voice of my supplication,'' say the same thing. God has no
body, he answers, and his ears are his power to hear: \latin{Deus ergo ubique
est, et ubique adest, audit, videt, efficit}, ``God therefore is everywhere,
and everywhere present; he hears, sees and acts.'' He is to be asked all the
same: \latin{Sed orandus a nobis est, ut secundum precem nostram adsit, audiat,
videat, efficiat. Naturae suae est ut audiat; sed fidei nostrae est ut
precemur audiri}. ``But he is to be prayed to by us, so that he may be
present, hear, see and act according to our prayer. It is of his nature that
he hears; but it is of our faith that we pray to be heard.'' He hears those
who revile him and those who speak in the secret of the heart; but faith does
its own office and asks to be heard, that he \latin{qui per naturam suam
audit, per orantis precem dignetur audire}, ``who hears by his nature, may
deign to hear through the prayer of one who prays'' (\work{Tractatus in Ps.}\
129.3, PL~9 cols.~719--720).

St John Chrysostom, on the same verses, takes the order of the words. The
psalmist first says ``I cried,'' and only then ``hear''; it is not possible,
Chrysostom says, simply to obtain from God unless one's own part has been
supplied, and intent prayer full of compunction has great power to persuade
God. In praying one should seek not only to receive what one asks but to make
the soul better (\work{Expositio in Ps.}\ 129, §1, PG~55).

Neither Hilary nor Chrysostom speaks of a captivity, a promise or a return.
Hilary answers from God's side, who hears by nature and is asked by faith;
Chrysostom from the side of the one who prays, whose own part comes first.
Their answers bear on Jerome's question. Because the same Mass sings the
command \latin{invocábitis me} and the cry \latin{De profúndis}, those who cry
from the depths are the people to whom the command was given, and their cry
is the call it asks of them.

\subsection{St Augustine: whose captivity, and which deep}
\label{sec:captivity-augustine}

St Augustine names the captivity. Of the Introit's verse, ``Thou hast turned
away the captivity of Jacob,'' he asks to whom it is said, and answers, ``To
Christ.'' He says which captivity Christ turns away: ``Thou hast
turned away our captivity, not by setting us free from the barbarians, with
whom we had not met, but by setting us free from bad works, from our sins, by
which Satan held sway over us'' (\work{Enarr.\ in Ps.}\ 84.2). The verb is
past, and Augustine explains the tense from the psalm's own next lines. After
saying that God has turned away his anger, the psalmist prays, ``Convert us,
O God our saviour: and turn off thy anger from us'' (v.~5). Augustine has him
explain: ``These things I speak of as done, because I see them about to be
done: but because they are not yet done, I pray that they may come, which I
have already seen'' (84.4).

The Postcommunion prays in the same way, for what is already pledged and
not yet done: it asks for the effect of salvation.

St Robert Bellarmine reads the verse of God's eternal decree to end the
captivity, a deliverance that ``will be thoroughly accomplished in the
resurrection only,'' and takes Jacob to mean ``the whole human race''
(\work{Explanatio in Ps.}\ 84, on v.~2). Cassiodorus also reads the
captivity as that of sin and says the psalmist rejoices \latin{tanquam de
praeterito}, as of a thing already past (\work{Expositio Psalmorum}, on Ps
84:2); the
exposition printed under St Athanasius's name reads the verse of Christ's
redemption of many from spiritual captivity, with the words of Isaias that he
read at Nazareth (Lk 4:18--21; \work{Expositio in Ps.}\ 84, on v.~2).

Of the Gradual's psalm Augustine says, ``Let us then hear in this Psalm the
voice of the Martyrs,'' because St Paul quotes its twenty-second verse of the
persecuted Christians (\work{Enarr.\ in Ps.}\ 43.1). Of the Alleluia's psalm,
whose deep is ``this mortal life,'' he says that the one who cries from it is
``a sinner,'' and that his hope is in the next verse, ``For with Thee there is
propitiation'': ``And what is this propitiation, except sacrifice? And what is
sacrifice, save that which hath been offered for us?'' (\work{Enarr.\ in Ps.}\
129.1--3).

The Mass does not sing that fourth verse. After the third \latin{De profúndis}
the Secret offers \latin{hóstias \dots\ placatiónis}, offerings of
propitiation; its word is \latin{placatio}, not the psalm's
\latin{propitiatio}. The words differ, but Augustine says the propitiation in
that verse is a sacrifice, and the Secret offers one; its first petition,
that God absolve our sins, asks for what the sinner who has cried from the
deep hopes for.

\subsection{The Gradual's past tense}
\label{sec:captivity-gradual}

The Gradual sings two verses of thanks, \latin{Liberásti nos \dots\
confudísti}, and stops before the psalm turns to lament: ``But now thou hast
cast us off, and put us to shame'' (Ps 43:10). Who can give thanks for a
deliverance in a psalm that complains of being abandoned? The Fathers answer
in different ways.

St Augustine says that the verse is ``spoken of the future under the figure
of the past,'' and that the reason ``is that it is as certain as if it were
past.'' He finds proof in the next verse, which the Gradual also sings,
\latin{In Deo laudábimur}: ``Observe how he intermingles words expressive of
a future time'' (\work{Enarr.\ in Ps.}\ 43.6--7). Only after that does
Augustine ask ``what is our state at present?'' and answer it with the
lament of verse~10. Cassiodorus agrees, \latin{Confidentia fidei futura pro
praeteritis dicunt}, ``by the confidence of faith they speak of things to come
as past,'' and places the saints' deliverance at the second coming
(\work{Expositio Psalmorum}, on Ps 43:8).

St Ambrose and St Thomas Aquinas put the deliverance both now and later.
Ambrose says the enemies are confounded here, \latin{sed multo amplius
confundentur; cum viderint sanctos Domini resurgentes}, ``but they will be
confounded far more when they see the Lord's saints rising'' (\work{Enarr.\
in Ps.}\ 43.21--22). Aquinas says God saves \latin{vel in praesenti \dots\
vel in futuro}, in the present or in the future (\work{Super Ps.}\ 43 n.~3).

St Robert Bellarmine reads the verses as a memory that grounds trust: ``I
acknowledge my safety is not owing to my own strength, because, as often as I
have been rescued from any danger, you `have saved us from those that afflict
us,''' and God's people, ``whenever they shall be delivered from any
tribulation, will thank God for it'' (\work{Explanatio in Ps.}\ 43, on
vv.~8--9). St John Chrysostom, at the verse before, has the psalmist use his
weapons and cast the whole upon God, and tells his hearer, when he fights the
devil, to say that he trusts not in his arms, his strength or his
righteousness but in God's mercy. At ``thou hast saved us'' he gives the
psalmist's own words: there is no need to speak of the forefathers, since
``we ourselves'' hold many pledges of God's providence. And he says that
\latin{tota die} means the whole of life, so that we glory in God not only
when he helps but also when he forsakes (\work{Expositio in Ps.}\ 43, §6).

The Collect says of the family it prays for what Bellarmine's psalmist says
and what Chrysostom has his hearer say, that its trust rests on God and on
nothing of its own: it \latin{in sola spe grátiæ cæléstis innítitur}, leans
on the hope of heavenly grace alone. Its petition for protection is then the prayer
of the same people whose thanks the Gradual sings.

Whichever answer is taken, the Gradual's thanks are sung, like the exiles'
prayer, by people still afflicted, before the deliverance is complete.

\subsection{The Communion's answer}
\label{sec:captivity-communion}

The Communion's promise, \latin{quidquid orántes pétitis, crédite quia
accipiétis, et fiet vobis}, answers the Introit's \latin{invocábitis me, et
ego exáudiam vos}: what God promised the exiles who call on him, Christ
promises to whoever asks in faith. Jerome draws a comparable link at the
verse the antiphon omits, joining ``you shall seek me, and shall find me''
(Jer 29:13) to \latin{Petite, et accipietis} (Mt 7:7). He does not mention the
saying in Mark.

\subsection{The chants read in order by the medieval commentators}
\label{sec:captivity-commentators}

Three medieval commentators on the Mass read the Introit, Gradual, Offertory
and Communion as one sequence and give each a step. Rupert of Deutz says that
in these chants the Church gives thanks for the future salvation of the
remnant of Israel. The Introit's thoughts of peace are the
promise of the banquet of grace that Christ will keep with his brethren
according to the flesh, prefigured in Joseph, who received the brothers who
had sold him. In the Gradual, freed from the captivity of souls, they give
thanks; in the Offertory they say to the true Joseph what Joseph's brothers
said, \latin{Obsecramus ut obliviscaris sceleris fratrum}, ``we beseech thee
to forget the wickedness of thy brethren'' (Gen 50:17); and the Communion is
his answer, \latin{Nolite timere}, ``Fear not'' (Gen 50:19) (\work{De divinis
officiis} XII.23, PL~170 cols.~329--330). William Durandus gives the same
order and the same figure, and glosses the Introit's \latin{non
afflictionis}: \latin{ad hoc tendunt mala quae vobis infero: ut reconciliem
vos mihi}, ``the evils I bring on you tend to this, that I may reconcile you
to myself.'' In the Offertory, he says, the Church asks pardon as Joseph's
brothers did; and because Joseph, that is Christ, answered \latin{Nolite
timere, ego pascam vos}, ``Fear not, I will feed you,'' and \latin{quia
liberalis est Dominus, et cito dimittit}, ``because the Lord is generous and
quickly forgives,'' the Communion follows with its promise to those who ask
(\work{Rationale} VI.140). Sicard of Cremona reads the Introit as the Lord's
announcement that \latin{Israelis reliquiae salvabuntur}, the remnant of
Israel will be saved, the Gradual and the Offertory as their answer, and the
Communion as the Lord assuring them again of salvation, \latin{certificat de
salute} (\work{Mitrale} VI.23).

The exposition printed under St Athanasius's name says something similar of
the Introit's psalm: it shows the calling of the Gentiles and then the Jewish
people confessing after it (\work{Expositio in Ps.}\ 84, the opening
summary).

\begin{fourSenses}
\item[Literal.] Jeremias writes to the exiles in Babylon that God thinks
thoughts of peace toward them and will bring them back when they call on him;
Psalms~84, 43 and 129 thank God for a captivity turned and a deliverance
given, and cry to him from the depths; Jesus tells his disciples that prayer
made in faith receives what it asks.
\item[Allegorical.] Christ turns away the captivity of sin (Augustine,
Cassiodorus and the Athanasian exposition on Ps 84); the voice of the Gradual's
psalm is the voice of the martyrs (Augustine); the propitiation with God is
the sacrifice offered for us (Augustine on Ps 129:4).
\item[Moral.] Pray, because God, who could fulfil his promise without our
prayer, urges us to pray that we may receive it (Jerome on Jer 29:12); pray
because it is God's nature to hear and faith's part to ask to be heard
(Hilary); supply one's own part, crying before asking to be heard
(Chrysostom); confess oneself a sinner in the deep (Augustine); give thanks
in affliction for a deliverance not yet complete (Augustine).
\item[Anagogical.] We remain in the confusion of this age until we are found
worthy to receive the rest of the number seven; then God fulfils his promise
and brings us back \latin{ad locum nostrum Ecclesiam} (Jerome); the
deliverance is accomplished ``in the resurrection only'' (Bellarmine), when
the enemies are confounded far more at the rising of the saints (Ambrose);
and, as the medieval commentators read the chants, the remnant of Israel is
saved (Rupert, Sicard) and the true Joseph answers his brothers' plea,
``Fear not'' (Rupert, Durandus).
\end{fourSenses}
